% Corrected French for printed p.192; departures are append-only below.
\clearpage
\SourcePage{192}

\LayerTarget{tohoku:corollary:4.3.2:1}
\textsc{Corollaire 1.} \textit{Supposons que $A$ soit annulé par un
idéal bilatère $I$ de $O$ tel que $O/I$ soit un anneau semi-simple $K$. Alors
les deux suites spectrales du bicomplexe $\mathrm{Hom}_O(L,\Gamma \mathrm C(A))$ sont
triviales, et leurs termes initiaux s'identifient canoniquement à l'homologie
$\mathrm{Ext}_{\mathbf O}(X;\mathbf M,A)$ de ce bicomplexe, et on a plus
précisément des isomorphismes canoniques :}

\LayerTarget{tohoku:eq:4.3.4}
\[
\tag{4.3.4}
\mathrm{Ext}_{\mathbf O}^n(X;\mathbf M,A)=
\bigoplus_{p+q=n}
\mathrm{Hom}_K(\mathrm{Tor}_p^O(K,M),H^q(X,A)).
\]
\LayerTarget{tohoku:proof:corollary-4.3.2:1}
En effet, $\mathrm C(A)$ et par suite $\Gamma \mathrm C(A)$ sera aussi annulé par $I$, d'où on
conclut

\LayerTarget{tohoku:eq:4.3.5}
\[
\tag{4.3.5}
\mathrm{Hom}_O(L,\Gamma \mathrm C(A))=
\mathrm{Hom}_K(K\otimes_O L,\Gamma \mathrm C(A))
\]
\LayerTarget{tohoku:proof:kunneth-collapse-4.3.4}
et comme $K$ est semi-simple on peut appliquer la version la plus simple du
théorème de Künneth (due au fait que $\mathrm{Hom}_K$ est un bifoncteur exact) :
On trouve la trivialité des suites spectrales, plus la formule
$H(\mathrm{Hom}_K(K\otimes_O L,\Gamma \mathrm C(A)))=
\mathrm{Hom}_K(H(K\otimes_O L),H(\Gamma \mathrm C(A))),$
ce qui n'est autre que la formule (4.3.4).

\LayerTarget{tohoku:construction:relative-k-hyperhomology-spectral-sequences}
Supposons toujours que $A$ est annulé par un idéal bilatère $I$, mais ne
supposons rien sur $K=O/I$. On a (4.3.5), ce qui amène à considérer les suites
spectrales d'hyperhomologie du bifoncteur $\mathrm{Hom}_K$ relativement aux
complexes $K\otimes_O L$ et $\Gamma \mathrm C(A)$. Dans la première le terme initial

\LayerTarget{tohoku:formula:first-relative-k-initial-term}
$H^p(\mathrm{Ext}_K^q(K\otimes_O L,\Gamma \mathrm C(A)))$ est nul pour $q>0$, puisque
$K\otimes_O L$ est $K$-libre. Donc l'aboutissement s'identifie canoniquement à
$H(\mathrm{Hom}_K(K\otimes_O L,\Gamma \mathrm C(A)))$ qu'il s'agit précisément de
calculer. Dans la deuxième suite spectrale, le terme initial est
\LayerTarget{tohoku:formula:second-relative-k-initial-term}
$\displaystyle \bigoplus_{q'+q''=q}
\mathrm{Ext}_K^p(H^{q'}(K\otimes_O L),H^{q''}\Gamma \mathrm C(A)).$
Donc :

\LayerTarget{tohoku:corollary:4.3.2:2}
\textsc{Corollaire 2.} \textit{Supposons le $\mathbf O$-Module $A$
annulé par un idéal bilatère $I$ de $O$, posons $K=O/I$. Alors
$(\mathrm{Ext}_{\mathbf O}^n(X;\mathbf M,A))$ est aussi l'aboutissement d'un
troisième foncteur spectral, dont le terme initial est}

\LayerTarget{tohoku:eq:4.3.6}
\[
\tag{4.3.6}
\mathrm{III}_2^{pq}(A)=
\bigoplus_{q'+q''=q}
\mathrm{Ext}_K^p(\mathrm{Tor}_{q'}^O(K,M),H^{q''}(X,A)).
\]
\LayerTarget{tohoku:corollary:left-hereditary-k}
\textit{En particulier, si $K$ est un \textup{``}anneau héréditaire à gauche\textup{''}
\textup{[6, 1.5] (}par exemple un anneau principal\textup{)} on a une suite exacte canonique :}

\LayerTarget{tohoku:sequence:left-hereditary-k}
\[
\begin{split}
0\to{}&
\bigoplus_{p+q=n-1}
\mathrm{Ext}_K^1(\mathrm{Tor}_p^O(K,M),H^q(X,A))\\
\to{}&\mathrm{Ext}_{\mathbf O}^n(X;\mathbf M,A)
\to
\bigoplus_{p+q=n}
\mathrm{Hom}_K(\mathrm{Tor}_p^O(K,M),H^q(X,A))
\to0.
\end{split}
\]

\LayerTarget{tohoku:remark:4.3.2:1}
\textsc{Remarques.} 1. Supposons les $L_i$ isomorphes à des
$O^{n_i}$, alors (prop. 4.3.2) la première suite spectrale d'hyperhomologie du
foncteur $\mathrm{Hom}_O(M,-)$ pour le complexe $\Gamma \mathrm C(A)$ est celle du
théorème 4.2.1 dont le terme initial a donc l'interprétation remarquable
(4.2.5). Mais on fera bien attention que ce fait avait été déduit du cor. 3 à
prop. 4.2.4 et tient essentiellement à la formule (4.2.11) c'est-à-dire
\textit{au fait que le foncteur $\mathrm C$ \textup{``}résolvante canonique\textup{''} permute au
foncteur} $\mathbf{Hom}_{\mathbf O}(\mathbf L,-)$, si $L$ est un $O$-module isomorphe à un $O^n$
(grâce à quoi on est d'ailleurs dans les conditions du corollaire à la
proposition 2.5.2). C'est

\TohokuContinuation{193}
